\documentclass{resume_cn-v2}
\begin{document}
\pagenumbering{gobble}

\name{胡昊}
\contact{%
  +86 133 8752 1087 \;\textperiodcentered\;
  \href{mailto:felixhuhao@gmail.com}{felixhuhao@gmail.com} \;\textperiodcentered\;
  \href{https://github.com/felixhuhao}{github.com/felixhuhao}%
}

\vspace{-8pt}\cvsection{个人概述}
\cvsummary{\textbf{资深后端工程师，现任 AI 方向 CTO} —— 10 余年构建大规模、高可用分布式系统与数据基础设施，曾就职于 Google、Uber 与 Microsoft。}
\cvsummary{\textbf{构建面向生产的 Agentic AI 系统} —— RAG、对话式数据分析、MCP 驱动的多智能体编排，以 harness engineering（黄金集评测、确定性护栏、上下文管理、人在回路安全）保障系统质量。}
\cvsummary{\textbf{实战型工程领导者} —— 带领小团队贯穿从架构到交付，采用规范驱动、审查把关的 AI 辅助工作流，配合质量门禁与自动化验证。}

\cvsection{技术能力}
\cvskill{大模型与智能体}{智能体编排、智能体治理、工具 / 函数调用、MCP 集成、多智能体路由、上下文工程、LangChain / LangGraph / DeepAgents、OpenAI / Claude / DeepSeek}
\cvskill{RAG 与检索}{RAG（多模态）、混合检索、嵌入、重排序、查询优化、Milvus、Qdrant}
\cvskill{评测与安全}{LangSmith、RAGAS、OpenEvals、LLM 裁判评测、护栏、人在回路、可观测性}
\cvskill{编程语言}{Java \& Go \textbf{(Google readability)}、Python、SQL、C\#、C++}
\cvskill{后端与分布式系统}{微服务、API 设计、gRPC / Protobuf、REST、Thrift、Spring AI、FastAPI、SSE}
\cvskill{数据系统}{数据建模、PostgreSQL、MySQL、Redis、BigQuery、Kafka、Flink、ClickHouse、MongoDB}
\cvskill{云与基础设施}{Google Cloud、AWS、Docker、Kubernetes、Linux、CI/CD}

\cvsection{工作经历}

\cventry{元宝创意}{2025.02 -- 至今}{首席技术官}{中国 武汉}
\begin{itemize}
  \item \textbf{架构并交付 ERP SaaS 平台的统一 AI 层}，覆盖运营、数据分析与知识，带领 4 人团队；建立共享的评测与安全实践：可观测性、黄金集与回归评测套件。
  \item \textbf{运营 Copilot}：基于 \textbf{DeepAgents} 构建智能体运行框架，搭载委托式子智能体（销售 / 库存 / 订单 / 采购），通过角色权限约束的 \textbf{MCP} 工具、沙箱化执行、主动监控、可审计工具追踪实现安全运营，并以 \textbf{Spring AI / Java~21} 服务提供人在回路审批（加密绑定、一次性使用）以保障高风险写操作。
  \item \textbf{数据分析智能体}：交付面向非技术用户的多轮自然语言转 SQL 分析系统，基于 ERP \textbf{ClickHouse} 数仓，由 \textbf{Qdrant} 语义层支撑，将提示上下文减少 \textbf{\textasciitilde73\%}；增加了 OLAP 意图路由、确定性 \textbf{SQLGlot} 护栏、只读执行、有界 SQL 修复、结果等价性回归评测及自动图表生成。
  \item \textbf{知识助手}：交付基于 \textbf{LangGraph} 编排的企业文档 RAG，支持多源数据接入与 \textbf{Milvus} 混合检索；根据查询意图选择精准 / 均衡 / 广域检索策略，结合查询扩展、HyDE 与 RRF 重排序，强制引用溯源、无证据不答及基于角色的访问控制；引入 LLM 裁判与关键切片综合评分，配合查询级可观测性。
\end{itemize}

\cventry{自由职业}{2023.03 -- 2025.01}{独立软件工程师与顾问}{加拿大多伦多 / 中国 武汉}
\begin{itemize}
  \item \textbf{构建并部署云原生旅行指南 SaaS}，支持 QR 扫码 POI 音频导览、匿名评分与播放分析，以 Next.js + Go 驱动，Postgres 存储 POI 与评分，音频走 S3/CDN 分发，Docker 容器化部署。
  \item \textbf{为客户交付 RAG / 多智能体原型与安全护栏方案}，覆盖多智能体编排、对话式检索等场景。
\end{itemize}

\cventry{Uber}{2022.05 -- 2023.03}{高级软件工程师}{加拿大多伦多}
\begin{itemize}
  \item \textbf{构建 Uber 首个 CPG（快消品）广告功能 —— 店铺横幅广告}：作为 3 人团队的后端负责人，构建基于 Go/gRPC + DocStore 的服务，连接 UberEats Storefront 与 Ads 平台，使广告主可在各大零售连锁中按预算投放商品/品牌广告；\textbf{首季度即达到 10--30K QPS，p99 \textasciitilde140\,ms}。
  \item \textbf{重新设计 CPG 广告计划的数据模型}：将餐厅时代的单点商品 ID 替换为 GTIN 索引的多门店模型 —— 使全国性广告计划成为可能，消除数百万条重复记录，形成 CPG 广告数据基础。
  \item \textbf{推出店铺菜单赞助商品}：将第三方广告供应商接入 CPG 广告主干，赞助商品在所在子菜单中优先展示，通过 Kafka + Flink 实时推送曝光/点击事件用于计费，\textbf{即便接入外部供应商，p95 仍维持在 \textasciitilde180\,ms}，配合 Prometheus / Grafana 监控与依赖告警。
  \item \textbf{主导 CPG 广告分阶段生产上线}，并担任广告团队开发代表 —— 协调产品、数据科学、销售、DevOps、合作团队及外部广告供应商。
\end{itemize}

\cventry{Google}{2018.02 -- 2021.02}{软件工程师 —— Cloud SQL（数据集成 / 灾备）}{美国 森尼韦尔}
\begin{itemize}
  \item \textbf{将大规模实例集群的备份健康度从 99.92\% 提升至 99.97\%}（支撑 99.95\% SLA）：独立负责并加固备份验证流水线（快照挂载、数据库启动、校验和比对，在共享 GCE 实例池上运行），解决包括资源耗尽、启动失败、误报验证失败在内的主要故障类别。
  \item \textbf{实现 \href{https://cloud.google.com/sql/docs/mysql/backup-recovery/manage-standard-backups\#point-in-time_recovery}{PITR} 安全的事务日志保留}：构建备份 PD 快照清理机制，精简非必要的通用日志数据，保留恢复关键日志，实现通过最近备份结合日志回放的任意时间点恢复。
  \item \textbf{重新设计备份保留策略}：将按固定数量淘汰（优先删除最旧备份、不区分健康状态）替换为始终保留经流水线验证为健康的最新备份，确保每个实例都拥有可恢复的还原点。
  \item \textbf{主导 MySQL 高可用从流式主备复制向 Persistent Disk 同步复制的迁移}：构建贯穿控制面与数据面的 Java 工作流，自动转换旧版流式主备配对。
  \item \textbf{为 MySQL 与 PostgreSQL 的 \href{https://cloud.google.com/sql/docs/postgres/import-export/import-export-csv}{CSV 导入/导出增加可定制格式}}：扩展后端 REST API，用 Go 重构单次导入/导出操作，打通 CSV 数据流转闭环。
  \item \textbf{将 Cloud SQL PostgreSQL CSV 导入/导出功能交付至 GA}：加固导入/导出路径，防范注入攻击，采用按需最小权限数据库访问，并通过 Cloud SQL IAM 管控对 Google Cloud Storage 的访问。
  \item \textbf{作为数据面 on-call 工程师维护 99.95\% SLA} —— 处理 Java 后端与 Go 实例代理上的发版阻塞及客户投诉事故；指导初级工程师。
\end{itemize}

\cventry{Microsoft}{2014.01 -- 2017.10}{软件开发工程师}{美国 山景城}
\begin{itemize}
  \item \textbf{从零开始联合构建 People I Communicate With(PICW)}，3 人团队：以 C\# 开发 Exchange 基础信号采集组件，采集常用联系人与邮件交互信号，服务于 Contact Safety 与 Exchange 反垃圾邮件系统，支持安全等级分类、垃圾规则处理与垃圾 / 钓鱼邮件检测等功能。
  \item \textbf{独立负责 PICW}，扩展其信号覆盖面并持续丰富信号类型，主导跨团队集成，对接 Exchange 反垃圾邮件、Contact Safety 与 People Relevance 下游系统；其信号成为 People Relevance 的核心输入，驱动 Graph People API 的推荐、排序与轻量级 CRM 等功能。
  \item \textbf{为 People Relevance 构建监控基础设施}：日志看板、实时告警与性能计数器，为团队提供 Graph People API 服务的运维可观测性，提升问题定位与 on-call 响应效率。
  \item \textbf{主导 Hotmail / Outlook.com 向 Exchange 迁移的账户筛选工具}：基于 C\# 构建分布式功能画像服务，通过 SQL 大规模扫描 22 亿账户基数的使用特征，发现 40+ 项 Exchange 功能缺口，识别 4 亿+ 可迁移账户，并在缺口各轮修复后重新画像以支持逐批迁移。
  \item \textbf{参与 Exchange on-call 轮值}，覆盖 Hotmail / Outlook 服务，处理发版阻塞与客户投诉事故，快速排查并高质量修复，保障 4 亿+ 用户邮箱服务的在线稳定性。
\end{itemize}

\cventry{Epic Systems}{2012.06 -- 2013.11}{软件开发工程师}{美国 麦迪逊}
\begin{itemize}
  \item \textbf{构建并扩展 EMC2}：Epic 内部软件配置管理平台(VB.NET on InterSystems Caché)，涵盖版本控制、依赖管理与自动化构建 / 部署 —— 一套内部 CI/CD 系统。
  \item \textbf{开发客户版本发布与升级工具}（控制台与 Web），通过安全通信框架自动化版本检查与更新推送，确保医疗系统运行于当前受支持的版本。
  \item \textbf{扩展 Trial Dev Environment}，支持安全的高风险变更集成，并解决构建 / 部署相关升级问题。
\end{itemize}

\cventry{HCR ManorCare}{2011.01 -- 2012.04}{数据仓库开发工程师}{美国 托莱多}
\begin{itemize}
  \item \textbf{构建 SSIS ETL 管道}，为拥有 500+ 医疗机构的全国性医疗集团驱动企业数据仓库：从多数据源抽取、验证、转换并加载至一致的维度 / 事实表模型。
  \item \textbf{开发 SSRS 报表与看板}（交叉表、钻取、参数化、级联），支持财务与运营商业智能。
  \item \textbf{创建 SQL Server 对象}（存储过程、索引视图、UDF），支撑报表逻辑与 ETL 验证。
\end{itemize}

\cvsection{教育背景}
\cvedu{Bowling Green State University}{2010.08 -- 2012.05}{计算机科学硕士}{GPA 4.0/4.0 \;\textperiodcentered\; 美国 鲍灵格林}
\cvedu{武汉大学}{2006.08 -- 2010.06}{软件工程学士}{GPA 3.5/4.0 \;\textperiodcentered\; 中国 武汉}

\end{document}
